\section{Galilean Relativity}

Imagine dropping a ball off a tower such as the Leaning Tower of Pisa,
and the tower is spinning with the Earth.
Assume we have two different observers, named Alice and Bob.

\begin{figure}
    \resizebox{\linewidth}{!}{\input{assets/galilean-relativity.tikz}}
\end{figure}

Alice is standing on the ground, and despite the tower spinning with the Earth,
she sees the ball fall straight down.
Bob is floating above the Earth and sees the ball curve as it falls.
Neither of them are more "right" than the other.
The difference arises
because they are measuring from different \textbf{frames of reference}.

Now imagine a third observer Charlie.
Charlie is spinning around in a circle.
What Charlie observes doesn't matter because Charlie is crazy.
They are undergoing centripetal acceleration
so anything they measure is going to give crazy results.

\textbf{Galilean relativity} says that the laws of physics are the same
for all \textbf{inertial} frames of reference,
i.e., frames of reference with no relative acceleration.
